\documentclass[10pt,a4paper]{amsart}
\usepackage[margin=19mm]{geometry}
\usepackage{amsmath,amssymb,mathtools}
\usepackage[scr=rsfs,cal=euler]{mathalfa}
\usepackage{xfrac}
\usepackage[unicode,hidelinks,bookmarks=false]{hyperref}
\newcommand{\ie}{i.e.\ }
\newcommand{\cf}{cf.\ }
\newcommand{\resp}{respectively\ }
\newcommand{\Subsection}{\subsection}
\providecommand{\nobreakdash}{\nobreak\hbox{-}}
\setlength{\emergencystretch}{2em}
\multlinegap0pt
\usepackage{bm}
\usepackage{bbm}
\usepackage{dsfont}
\usepackage{xspace}
\usepackage{cancel}
\usepackage{mathtools}
\usepackage{amsthm}
\makeatletter
\@ifundefined{theorem}{\newtheorem{theorem}{Theorem}}{}
\makeatother
\title[Calibration Meets Reality: Making Machine Learning Predictio]{Calibration Meets Reality: Making Machine Learning Predictions Trustworthy}
\author{Kristina P. Sinaga, Arjun S. Nair}
\date{Source: arXiv:2509.23665v1 (CC BY 4.0)}
\begin{document}
\maketitle

\noindent\emph{Source and licence.} Calibration Meets Reality: Making Machine Learning Predictions Trustworthy. Version arXiv:2509.23665v1 (CC BY 4.0), distributed under the Creative Commons Attribution 4.0 International licence, as recorded in the arXiv licence metadata for this exact version and shown on the abstract page https://arxiv.org/abs/2509.23665v1. Full rights notice: licenses/arxiv-2509.23665-CC-BY-4.0.txt.\par\smallskip

\noindent\textit{Editorial scope.} Counted unit: the Theorem whose source label is \texttt{None} in the article (source0.tex lines 340--342, source label \texttt{None}), delivered together with the article's own complete original proof of it (source0.tex lines 344--396, one \texttt{proof} environment of 53 lines). No mathematics is written, repaired or re-derived: statement and proof are the article's own text line for line. Required context retained uncounted: none. Every definition, display and label of those lines is retained unchanged. Omitted from this package and not redistributed: 1--339, 343--343, 397--1185. Nothing inside a retained excerpt is omitted or altered: every retained body is delivered exactly as the article's lines are written, so no omission receipt of any kind accompanies this package. Citations: the retained excerpts carry no citation command at all, so no bibliography entry of the article is transcribed and no reference is redistributed. Reference adaptation: none was needed and none was made; every reference inside the delivered unit resolves inside it, so no unresolved reference or citation is produced. Printed slips kept literal: no printed word, symbol, hypothesis or number of the counted statement or of its proof was altered. Layout and class: the article's own class is \texttt{article} with packages including PRIMEarxiv, algorithm2e, amsmath, amssymb, amsfonts, amsthm, algorithm, algorithmic, bm, booktabs, cases, caption, cite, color; the wrapper is the lane's pinned \texttt{amsart} editorial wrapper at 10pt on A4, which restates the theorem-like environments the retained text needs and adds the article's own macro definitions that the retained text needs (none), with extra wrapper packages (bm, bbm, dsfont, xspace, cancel, mathtools). The delivered numbering is the unit's own. Assets: no retained span contains a figure environment, a graphics-inclusion command, a PDF-page inclusion, a table or a TikZ picture; no image file is copied into this package, the package declares no asset dependency and no image file is redistributed with it. Licence: version arXiv:2509.23665v1 of this article is distributed under the Creative Commons Attribution 4.0 International licence, as recorded in the arXiv licence metadata for this exact version and shown on the abstract page https://arxiv.org/abs/2509.23665v1. A full rights notice is delivered at licenses/arxiv-2509.23665-CC-BY-4.0.txt. Characters: every retained excerpt is pure ASCII, so the delivered engine, XeLaTeX with its default Latin Modern text font, reports no missing character.
\begin{theorem}[Uniqueness of Isotonic Regression Solution]
The isotonic regression problem has a unique solution $g^*$ that is a right-continuous step function with at most $n$ steps.
\end{theorem}

\begin{proof}
We establish the uniqueness and structure of the isotonic regression solution through a comprehensive analysis of the optimization problem's properties.

\textbf{Step 1: Problem Formulation and Constraint Set.}
The isotonic regression problem seeks to minimize:
$$J(g) = \sum_{i=1}^{n} (y_i - g(s_i))^2$$
subject to the constraint that $g: \mathbb{R} \rightarrow \mathbb{R}$ is non-decreasing, i.e., $s \leq t \Rightarrow g(s) \leq g(t)$.

Let $\mathcal{G}$ denote the set of all non-decreasing functions on $\mathbb{R}$. This constraint set is convex: for any $g_1, g_2 \in \mathcal{G}$ and $\lambda \in [0,1]$, the function $\lambda g_1 + (1-\lambda) g_2$ is also non-decreasing.

\textbf{Step 2: Strict Convexity Analysis.}
The objective function $J(g)$ is strictly convex in $g$. To see this, consider any two distinct functions $g_1, g_2 \in \mathcal{G}$ with $g_1 \neq g_2$. For any $\lambda \in (0,1)$:

\begin{align}
J(\lambda g_1 + (1-\lambda) g_2) &= \sum_{i=1}^{n} (y_i - \lambda g_1(s_i) - (1-\lambda) g_2(s_i))^2 \\
&= \sum_{i=1}^{n} (\lambda(y_i - g_1(s_i)) + (1-\lambda)(y_i - g_2(s_i)))^2
\end{align}

By the strict convexity of the squared function $x \mapsto x^2$, we have:
$$J(\lambda g_1 + (1-\lambda) g_2) < \lambda J(g_1) + (1-\lambda) J(g_2)$$
provided that $(y_i - g_1(s_i)) \neq (y_i - g_2(s_i))$ for at least one index $i$, which holds when $g_1 \neq g_2$.

\textbf{Step 3: Existence and Uniqueness via Optimization Theory.}
Since $\mathcal{G}$ is a convex set and $J(g)$ is strictly convex, the optimization problem has at most one solution. Existence is guaranteed by considering the restriction of $g$ to the finite set $\{s_1, s_2, \ldots, s_n\}$ and applying compactness arguments.

More precisely, we can restrict attention to functions $g$ such that $\min_i y_i \leq g(s_i) \leq \max_i y_i$ for all $i$, as any function outside this range cannot be optimal. This restriction creates a compact feasible region in the finite-dimensional space $\mathbb{R}^n$, ensuring existence of a minimizer.

\textbf{Step 4: Characterization as a Step Function.}
We now establish that the optimal solution $g^*$ is a step function. The key insight is that $g^*$ need only be defined on the set $\{s_1, s_2, \ldots, s_n\}$ to solve our problem, and its extension to all of $\mathbb{R}$ can be chosen optimally.

Let $s_{(1)} \leq s_{(2)} \leq \ldots \leq s_{(k)}$ denote the distinct values among $\{s_1, s_2, \ldots, s_n\}$, where $k \leq n$. For any non-decreasing function $g$, we must have:
$$g(s_{(1)}) \leq g(s_{(2)}) \leq \ldots \leq g(s_{(k)})$$

The optimal choice of $g^*$ on each equivalence class of equal $s_i$ values is determined by minimizing the sum of squared residuals within that class, which yields the average of corresponding $y_i$ values.

\textbf{Step 5: Right-Continuity and Step Structure.}
For $s \in (s_{(j)}, s_{(j+1)})$, the value of $g^*(s)$ does not affect the objective function. The optimal choice is $g^*(s) = g^*(s_{(j)})$ to maintain the non-decreasing property with minimal "jumps."

At each $s_{(j)}$, we define $g^*(s_{(j)})$ as the right limit to ensure right-continuity:
$$g^*(s_{(j)}) = \lim_{s \to s_{(j)}^+} g^*(s)$$

This construction yields a right-continuous step function with at most $k \leq n$ steps, where each step occurs at some $s_{(j)}$.

\textbf{Step 6: Pool Adjacent Violators Characterization.}
The solution $g^*$ can be characterized through the Pool Adjacent Violators (PAV) algorithm. The optimal values satisfy:
$$g^*(s_{(j)}) = \frac{\sum_{i: s_i \in [s_{(j_1)}, s_{(j_2)}]} y_i}{\sum_{i: s_i \in [s_{(j_1)}, s_{(j_2)}]} 1}$$
for appropriate intervals $[s_{(j_1)}, s_{(j_2)}]$ determined by the PAV merging process.

\textbf{Step 7: Uniqueness Conclusion.}
The combination of strict convexity of the objective function and convexity of the constraint set ensures that the solution is unique. The step function structure with at most $n$ steps follows from the finite sample nature of the problem and the optimality conditions derived above.

Therefore, the isotonic regression problem has a unique solution $g^*$ that is a right-continuous step function with at most $n$ steps, completing the proof.
\end{proof}

\end{document}
